\section{Agentes que se automejoran: transformar la generación única en un sistema de ciclo cerrado}

Post-entrenamiento, el modelo aún puede verse como una función de entrada a salida; los agentes permiten que el modelo perciba, decida, actúe y actualice su estado repetidamente en un entorno. Lo crucial no es usar el nombre "agente", sino si el sistema posee objetivos, observación, espacio de acción, memoria, permisos de herramientas y condición de terminación.

\subsection{Definición operativa del agente}

Supongamos que el estado del entorno es $s_t$, la observación visible para el modelo es $o_t$, la memoria interna es $m_t$ y la estrategia de acción es
\[
a_t\sim\pi_{\theta}(a\mid o_t,m_t,g),
\]
donde $g$ es el objetivo. El entorno devuelve $o_{t+1}$ y la actualización de memoria es $m_{t+1}=U(m_t,o_{t+1},a_t)$. La capacidad del agente proviene del ciclo cerrado formado por la estrategia, las herramientas, la actualización de estado y la retroalimentación del entorno, no está determinada solo por los parámetros del modelo.

\lecturefigure{slide-102.jpg}{Componentes del agente de IA: objetivo, decisión, herramientas, memoria y retroalimentación}{CS25 V6 official deck, p.~102}

\lecturefigure{slide-103.jpg}{Ciclo del agente: percibir, razonar, actuar, observar, reflexionar y repetir}{CS25 V6 official deck, p.~103}

\teachervoice{00:43:12--00:44:05, el docente comprime la diferencia entre agentes y modelos de una sola vez en un ciclo. Nota de clase: los prompts multironda sin observación o actualización de estado no necesariamente constituyen una interacción ambiental real.}

\begin{importantbox}{Campos mínimos del contrato del agente}
El objetivo $g$, la observación visible, las acciones/herramientas permitidas, la política de memoria, el presupuesto, la condición de parada y los límites de recuperación de errores y permisos deben definirse explícitamente. Si falta cualquiera de estos elementos, la "autonomía" se convertirá en una promesa ambigua no auditable.
\end{importantbox}

\subsection{Automejora: modificar salida, memoria o estrategia}

La "automejora" puede referirse a tres niveles distintos: refinamiento de la respuesta actual; escribir resúmenes de fallos en la memoria episódica; o actualizar parámetros mediante entrenamiento. La diapositiva 104 propone el marco general, mientras que las diapositivas 105--106 muestran el refinamiento y la reflexión por separado. Los dos primeros alteran principalmente el contexto actual o el estado externo, lo cual no equivale a que el modelo haya aprendido permanentemente nuevas capacidades.

\lecturefigure{slide-104.jpg}{Automejora: generación, retroalimentación, revisión y acumulación de experiencia}{CS25 V6 official deck, p.~104}

\lecturefigure{slide-105.jpg}{Refinamiento: crítica e iteración revisada del mismo output}{CS25 V6 official deck, p.~105}

\lecturefigure{slide-106.jpg}{Reflexión: resumir fallos en memoria reutilizable para tareas posteriores}{CS25 V6 official deck, p.~106}

Si la candidata $k$-ésima es $y^{(k)}$, el crítico genera retroalimentación $c^{(k)}$, y el refinamiento puede escribirse como
\[
y^{(k+1)}\sim p_{\theta}(y\mid x,y^{(k)},c^{(k)}).
\]
Solo cuando los errores del crítico y del generador no están completamente correlacionados, la retroalimentación puede localizar problemas reparables y el presupuesto permite múltiples iteraciones, es posible que el ciclo mejore el rendimiento. De lo contrario, el sistema repetirá la reescritura de redacción o amplificará el mismo error.

\subsection{ReAct: alternar razonamiento y acción}

ReAct alterna pensamiento, acción y observación en la trayectoria. La acción llama a búsqueda, bases de datos, código u otras APIs; la observación devuelve los resultados reales al modelo. Tiene más anclaje ambiental que el CoT puro, pero aumenta errores de selección de herramientas, parámetros incorrectos, violación de permisos, inyección de prompts y estados externos no reproducibles.

\lecturefigure{slide-107.jpg}{Ciclo cerrado alternado de razonamiento, acción de herramienta y observación en ReAct}{CS25 V6 official deck, p.~107}

\begin{lstlisting}[language=Python,caption={Ciclo mínimo del agente con permisos y condición de terminación}]
state = initialize(goal, memory)
for step in range(max_steps):
    proposal = model.plan(state)
    if proposal.action == "finish":
        return proposal.answer
    authorize(proposal.action, policy=tool_policy)
    observation = tools.execute(proposal.action)
    state = update_state(state, proposal, observation)
raise BudgetExceeded(state)
\end{lstlisting}

\teachervoice{00:44:05--00:45:43, el docente considera que el refinamiento, la reflexión y ReAct consisten en insertar retroalimentación en diferentes ubicaciones: después del output, memoria entre tareas o después de una acción ambiental. Experiencia práctica: a mayor longitud del ciclo, más se requiere ingeniería para el estado, permisos y terminación.}

\begin{warningbox}{Más iteraciones no garantizan fiabilidad}
Si el crítico y el generador comparten puntos ciegos, la reflexión resumirá errores en memoria más persuasiva; si la observación de las herramientas está contaminada por inyección de prompts, la siguiente razonación continuará desde un estado erróneo. La evaluación del agente debe medir tasa de éxito, costo, capacidad de recuperación y efectos secundarios, no solo la calidad final del texto.
\end{warningbox}

La evaluación del agente debe tomar la trayectoria en lugar de la respuesta final como unidad mínima de auditoría. Se necesita registrar cada llamada a herramienta, parámetros, origen de observación, cambios de estado y causa de parada, luego estadizar éxito de tarea, pasos inválidos, acciones peligrosas, tasa de recuperación, costo token/latencia y violaciones de permisos. Para herramientas irreversibles, también se debe distinguir simulación, ejecución previa a producción y ejecución real. Si el benchmark solo otorga recompensa final, el sistema podría tener éxito mediante atajos accidentales ocultando rutas inseguras; si solo se mira una trazada bonita, podría recompensar explicaciones extensas. El objetivo de un agente confiable es completar tareas de manera estable con costos controlables y detenerse en un estado seguro ante fallos.

\subsection{Resumen de la sección}

El agente transforma el modelo en una estrategia de ciclo cerrado, pero cada interfaz del ciclo introduce nuevas capacidades de observación y patrones de fallo. El refinamiento modifica la candidata actual, la reflexión cambia la experiencia externa y ReAct añade observación real. Para hablar de automejora, es necesario especificar si lo que se modifica es el output, memoria, estrategia de herramientas o parámetros.

